% ID: SEL-D0046-OB-E03
% ESTADO: validada. S047.
Sean $C$ «Antonio compra» y $F$ «la fruta está de oferta». Tenemos $P(C)=2/5$, $P(C^c)=3/5$, $P(F\mid C)=1/3$ y $P(F\mid C^c)=1/2$.

\textbf{a)} Por probabilidad total,
\[
P(F)=\tfrac25\cdot\tfrac13+\tfrac35\cdot\tfrac12=\boxed{\tfrac{13}{30}}.
\]
\textbf{b)} Sumamos los días de compra y los días sin compra en los que hay oferta. Los dos sucesos son disjuntos:
\[
P(C\cup F)=P(C)+P(C^c)P(F\mid C^c)=\tfrac25+\tfrac35\cdot\tfrac12=\boxed{\tfrac7{10}}.
\]
